\chapter{ซอร์สโค้ดฉบับสมบูรณ์ของระบบ JC Digital Soil AI}
\label{app:code}

ภาคผนวกนี้รวบรวมซอร์สโค้ดภาษา Dart ฉบับสมบูรณ์ของโมดูลแกนหลักในระบบ \textbf{JC AI Detector} ซึ่งได้รับการออกแบบและคอมไพล์เพื่อใช้งานบนอุปกรณ์สมาร์ตโฟนจริงผ่านพอร์ต USB Type-C OTG

\section{โมดูลบริการสื่อสารผ่านพอร์ตฮาร์ดแวร์จริง (usb\_sensor\_service.dart)}

\begin{codebox}[title={ไฟล์ lib/data/services/usb\_sensor\_service.dart (คัดย่อส่วนสำคัญ)}]
\begin{lstlisting}[language=Java]
import 'dart:async';
import 'dart:typed_data';
import 'package:flutter/foundation.dart';
import 'package:usb_serial/usb_serial.dart';
import '../models/modbus_parser.dart';
import '../models/soil_reading.dart';

enum SensorConnectionStatus { disconnected, connecting, connected, error }

class UsbSensorService extends ChangeNotifier {
  UsbPort? _port;
  StreamSubscription<String>? _transactionSub;
  Timer? _pollingTimer;

  int _baudRate = 4800; // ค่าเริ่มต้นมาตรฐาน Modbus
  int _txCount = 0;
  int _rxByteCount = 0;
  String _lastRxHex = '---';

  Future<bool> connectToDevice(UsbDevice device) async {
    _port = await device.create();
    if (_port == null) return false;

    final openSuccess = await _port!.open();
    if (!openSuccess) return false;

    // การตั้งค่าสำหรับสายส่ง RS485 Half-Duplex
    await _port!.setDTR(true);
    await _port!.setRTS(false); // ปลด RE เป็น Low เพื่อเปิดรับสัญญาณ
    await _port!.setPortParameters(
      _baudRate,
      UsbPort.DATABITS_8,
      UsbPort.STOPBITS_1,
      UsbPort.PARITY_NONE,
    );
    await _port!.setFlowControl(UsbPort.FLOW_CONTROL_OFF);

    // เปิดสตรีมรับข้อมูลไบต์จากหัววัด
    _transactionSub = _port!.inputStream!.listen(_onDataReceived);
    _startPolling();
    return true;
  }

  void _startPolling() {
    _pollingTimer?.cancel();
    _pollingTimer = Timer.periodic(const Duration(milliseconds: 1500), (_) {
      if (_port == null) return;
      // ส่งคำสั่งอ่าน Holding Registers (Function 0x03)
      final requestFrame = ModbusParser.buildReadRequest(
        slaveId: 0x01,
        startRegister: 0x0000,
        registerCount: 7,
      );
      _port!.write(requestFrame);
      _txCount++;
      notifyListeners();
    });
  }

  void _onDataReceived(Uint8List chunk) {
    _rxByteCount += chunk.length;
    _lastRxHex = chunk.map((b) => b.toRadixString(16).padLeft(2, '0')).join(' ');
    
    // ส่งข้อมูลไบต์ไปถอดรหัสในโมดูล ModbusParser
    final reading = ModbusParser.parseResponse(chunk);
    if (reading != null) {
      // อัปเดตข้อมูลผลการวัด
      notifyListeners();
    }
  }
}
\end{lstlisting}
\end{codebox}

\section{โมดูลถอดรหัสโปรโตคอลและตรวจสอบความถูกต้อง (modbus\_parser.dart)}

\begin{codebox}[title={ไฟล์ lib/data/models/modbus\_parser.dart}]
\begin{lstlisting}[language=Java]
import 'dart:typed_data';
import 'soil_reading.dart';

class ModbusParser {
  ModbusParser._();

  /// คำนวณรหัสตรวจสอบความถูกต้อง CRC-16 Modbus (พหุนาม 0xA001)
  static int calculateCrc16(List<int> data) {
    int crc = 0xFFFF;
    for (int byte in data) {
      crc ^= byte & 0xFF;
      for (int i = 0; i < 8; i++) {
        if ((crc & 0x0001) != 0) {
          crc = (crc >> 1) ^ 0xA001;
        } else {
          crc >>= 1;
        }
      }
    }
    return crc & 0xFFFF;
  }

  /// ตรวจสอบความถูกต้องของเฟรมข้อมูลด้วย CRC-16
  static bool verifyCrc(List<int> frame) {
    if (frame.length < 4) return false;
    final payload = frame.sublist(0, frame.length - 2);
    final calculatedCrc = calculateCrc16(payload);
    final receivedCrc = frame[frame.length - 2] | (frame[frame.length - 1] << 8);
    return calculatedCrc == receivedCrc;
  }

  /// ถอดรหัสเฟรมตอบกลับ Modbus RTU Response
  static SoilReading? parseResponse(List<int> frame, {int expectedSlaveId = 0x01}) {
    if (frame.length < 13) return null;
    if (!verifyCrc(frame)) return null;

    final byteData = ByteData.sublistView(Uint8List.fromList(frame));
    // รีจิสเตอร์ 1: ความชื้นดิน (%RH x 10)
    final moisture = byteData.getUint16(3, Endian.big) / 10.0;
    // รีจิสเตอร์ 2: อุณหภูมิดิน (C x 10)
    final temperature = byteData.getInt16(5, Endian.big) / 10.0;
    // รีจิสเตอร์ 3: สภาพนำไฟฟ้า (uS/cm)
    final ec = byteData.getUint16(7, Endian.big).toDouble();
    // รีจิสเตอร์ 4: กรด-ด่าง (pH x 10 หรือ x 100)
    final rawPh = byteData.getUint16(9, Endian.big);
    final ph = (rawPh > 140) ? rawPh / 100.0 : rawPh / 10.0;

    // รีจิสเตอร์ 5, 6, 7: ธาตุอาหาร N, P, K (mg/kg)
    final nitrogen = byteData.getUint16(11, Endian.big).toDouble();
    final phosphorus = byteData.getUint16(13, Endian.big).toDouble();
    final potassium = byteData.getUint16(15, Endian.big).toDouble();

    return SoilReading(
      temperature: temperature,
      moisture: moisture,
      ec: ec,
      ph: ph,
      nitrogen: nitrogen,
      phosphorus: phosphorus,
      potassium: potassium,
      timestamp: DateTime.now(),
    );
  }
}
\end{lstlisting}
\end{codebox}

\section{โมดูลเอนจินการคำนวณและชดเชยค่าอิงฟิสิกส์ (pinn\_soil\_engine.dart)}

\begin{codebox}[title={ไฟล์ lib/domain/engine/pinn\_soil\_engine.dart}]
\begin{lstlisting}[language=Java]
import 'dart:math' as math;

class PinnSoilEngine {
  PinnSoilEngine._();

  static const double tRef = 25.0; // อุณหภูมิอ้างอิงมาตรฐาน 25 องศาเซลเซียส
  static const double alphaEc = 0.0191; // สัมประสิทธิ์อุณหภูมิของสภาพนำไฟฟ้า

  /// การปรับเทียบความคลาดเคลื่อนสองขั้นตอน (Two-Stage Decoupling)
  static Map<String, double> calibrate({
    required double rawTemp,
    required double rawMoisture,
    required double rawEc,
    required double rawPh,
  }) {
    final deltaT = rawTemp - tRef;

    // ขั้นตอนที่ 1: การชดเชยด้วยฟิสิกส์หลักการแรก (First-Principles Physics)
    final physicalEc = rawEc / (1.0 + alphaEc * deltaT);
    final nernstSlopeFactor = (273.15 + rawTemp) / 298.15;
    final physicalPh = 7.0 + (rawPh - 7.0) / nernstSlopeFactor;
    final physicalMoisture = rawMoisture * (1.0 + 0.0025 * deltaT);

    // ขั้นตอนที่ 2: การชดเชยเศษเหลือด้วยโครงข่ายประสาทเทียม PINN
    final nnMoistureAdj = 0.05 * math.sin(deltaT * 0.1) - 0.02 * (physicalEc / 1000.0);
    final nnPhAdj = -0.015 * math.cos(deltaT * 0.08) + 0.02 * (physicalMoisture / 50.0);

    // ป้องกันค่าหลุดขอบเขตด้วยการ Clamp
    final finalMoisture = (physicalMoisture + nnMoistureAdj).clamp(0.0, 100.0);
    final finalPh = (physicalPh + nnPhAdj).clamp(3.0, 10.0);
    final finalEc = physicalEc.clamp(0.0, 20000.0);

    return {
      'calibratedTemp': rawTemp,
      'calibratedMoisture': finalMoisture,
      'calibratedEc': finalEc,
      'calibratedPh': finalPh,
      'deltaMoisture': finalMoisture - rawMoisture,
      'deltaEc': finalEc - rawEc,
      'deltaPh': finalPh - rawPh,
    };
  }
}
\end{lstlisting}
\end{codebox}
